\documentclass[12pt]{article}

\usepackage{xcolor}
\renewcommand\thefootnote{\textcolor{red}{\arabic{footnote}}}
\usepackage[margin=1in]{geometry}
\usepackage{fancyhdr,graphicx, parskip}
\usepackage[spanish]{babel}
\usepackage{amsmath,amssymb,amsfonts,latexsym,cancel}\usepackage{rawfonts}
\usepackage{pictexwd}


\renewcommand{\headrulewidth}{0pt}
\fancyhead[L]{\includegraphics[width=3cm]{logo_uv.eps}}




\begin{document}
	
	\title{\textbf{Soluciones Tarea Obligatoria II\\ M\'etodos Matem\'aticos F\'isica II }}
	\author{\textit{$^{1}$Vicente Herrera, $^{2}$Marcelo \'Ordenes, $^{3}$Felipe Ortiz, $^{4}$Fabi\'an Trigo.}\\
		\small{\textit{$^{1,2,3,4}$Estudiantes de 3er. a\~{n}o de la Licenciatura en F\'isica}}\\
		\small {\textit{Universidad de Valpara\'iso, Facultad de Ciencias, CHILE.}}\\
		\small{\texttt{$^{1}$\textit{vicente.herrera@alumnos.uv.cl}}}\\
		\small{\texttt{$^{2}$\textit{marcelo.ordenes@alumnos.uv.cl}}}\\
		\small{\texttt{$^{3}$\textit{felipe.ortiza@alumnos.uv.cl}}}\\
		\small{\texttt{$^{4}$\textit{fabian.trigo@alumnos.uv.cl}}}
	}
	\date{\small{\today}}
	
	\maketitle
	\thispagestyle{fancy}

\rule{112mm}{0.1mm}\\
\textbf{Problema I: Campo El\'ectrico y plano infinito cargado}\\
\rule{112mm}{0.1mm}
\begin{enumerate}
	\item [A)] Del enunciado tenemos que el plano es infinito y est\'a cargado en toda su extensi\'on, es decir, $\sigma_{0} = cte.$ Adem\'as, esta distribibuci\'on de carga coincide con el plano xy. Por otro lado, el campo el\'ectrico vectorialmente es perpendicular al plano cargado y su magnitud es $|\vec{E}|= \frac{\sigma_{0}}{2 \epsilon_{0}}$. Se nos pide demostrar para el punto $\vec{R} = (0,0,z_{0})$ que la componente $z$ del campo el\'ectrico se puede escribir como:
	\[ E_{z} = \vec{E} \cdot \hat{k} = |\vec{E}| = ALGO \cdot \int_{-\infty}^{\infty} \frac{dx dy}{(x^{2}+y^{2}+z_{0}^{2})^{\frac{3}{2}}} \]
	Y determinar ese factor ALGO.
	
	\item [B)]
	\item [C)]
	\item [D)] 
\end{enumerate}


\rule{112mm}{0.1mm}\\
\textbf{Problema II: MoB y funciones singulares en el origen}\\
\rule{112mm}{0.1mm}\\
\begin{enumerate}
	\item [A)]
	\item [B)]
	\item [C)]
	\item [D)]
	\item [E)]
\end{enumerate}


\rule{112mm}{0.1mm}\\
\textbf{Problema III: Integraciones varias}\\
\rule{112mm}{0.1mm}
\begin{enumerate}
	\item [A)]
	\item [B)]
	\item [C)]
\end{enumerate}



\end{document}